\section{Выводы}

В ходе выполнения лабораторной работы по курсу «Дискретный анализ» была реализована система эффективного поиска множества паттернов в заранее известном тексте на основе \textbf{суффиксного массива}.

Основные результаты работы:
\begin{itemize}
    \item Реализован препроцессинг входного текста с использованием \textbf{алгоритма Манбера-Майерса}. За счёт применения итеративного удвоения длины префиксов и стабильной сортировки подсчётом удалось достичь временной сложности построения индекса $O(N \log N)$, где $N$ — длина текста.
    \item Спроектирована архитектура \textbf{быстрого поиска}, основанная на бинарном поиске по лексикографически отсортированным суффиксам. Это обеспечивает сложность поиска $O(M \log N)$ для каждого запроса (где $M$ — длина паттерна), что делает алгоритм высокоэффективным при многократном поиске различных образцов.
    \item Реализован механизм \textbf{потоковой обработки запросов}: после однократной индексации текста программа способна принимать на вход непрерывный поток паттернов. Использование модифицированных функций бинарного поиска (\texttt{lower\_bound} и \texttt{upper\_bound}) с кастомными компараторами позволило мгновенно локализовывать границы совпадений и восстанавливать исходные индексы вхождений.
\end{itemize}

\pagebreak